\cotefile{79}
\scanpage{250}

\unsure{demment|évidemment} surjectif, dont le noyau est le groupe

\[
\unsure{\mathfrak{z}|\mathfrak{Z}}=\deleted{\unsure{G_0}}G_{(s_0,t_0)}=\{g\in G\mid g(s_0)=s_0,\ g(t_0)=t_0\}\tag*{(91)}
\]

% slanted
\margin{NB. On a $\mathfrak{z}=G_0(s_0,t_0)$}

d'où on a une suite exacte

\[\begin{aligned}
1&\unsure{-|\longrightarrow}\mathfrak{z}\unsure{-|\longrightarrow}T_G\longrightarrow Z\times Z\unsure{-|\longrightarrow}1\\
&\hspace{5em}g\longmapsto(\delta_0(g),\delta_1(g))
\end{aligned}\tag*{(92)}\]

On pose

\[
\left\{\begin{aligned}
T_{0\deleted{G}}&=\operatorname{Ker}\delta_0=\deleted{\unsure{G_{s_0}}\unsure{\cap T_G}=}\{g\in G\mid gs_0=s_0,\,gt_0\in Zt_0\}\\
T_{1\deleted{G}}&=\operatorname{Ker}\delta_1=\{g\in G\mid gs_0\in Zs_0,\,gt_0=t_0\}
\end{aligned}\right.\tag*{(93)}
\]

Notons qu'on a

\[
\left\{\begin{aligned}
&\deleted{\unsure{T_G}}\mathfrak{z}\subset T_{0\deleted{G}},T_{1\deleted{G}}\subset\deleted{\unsure{G_0}}T_G,\quad Z\subset T_G\\
&T_G=T_{0\deleted{G}}\cdot T_{1\deleted{G}}=T_{0\deleted{G}}\cdot Z=\deleted{Z}T_{1\deleted{G}}\cdot\deleted{Z}\\
&T_{0\deleted{G}}\cap T_{1\deleted{G}}=T_{0\deleted{G}}\cap Z=Z\cap T_{1\deleted{G}}=\mathfrak{z}
\end{aligned}\right.\tag*{(94)}
\]

et

\[
\left\{\begin{aligned}
&\sigma(T_G)=T_G\\
&\sigma(T_{0\deleted{G}})=T_{1\deleted{G}},\quad \sigma(T_{1\deleted{G}})=T_{0\deleted{G}}\\
&\text{et bien sûr}\qquad \sigma(Z)=Z\quad (Z\text{ central }?)
\end{aligned}\right.\ .\tag*{(95)}
\]

On a aussi

\[\begin{aligned}
T_G\cap H=T_H&\subset\operatorname{Ker}(\delta_0.\delta_1:T_G\longrightarrow Z)\\
&\hspace{9em}\deleted{\unsure{g}}t\longmapsto\delta_0(t)\delta_1(t)
\end{aligned}\tag*{(96)}\]

et \deleted{\unsure{sur}} $\delta_H:T_H\xrightarrow{\sim}Z$ est \unsure{égale à la|opposée à la} restriction de $\unsure{\delta_0}$ à $T_H$

\deleted{Introduisons maintenant}

\[\deleted{T_G\ \text{normalise }U_0\text{ et centralise }T_H}\tag*{\deleted{(97)}}\]

\[\deleted{\operatorname{Ker}(\delta_0\delta_1:T_G\longrightarrow Z)=T_H}\tag*{\deleted{(97)}}\]

\deleted{Comme \unsure{$G$} $T_{G\unsure{0,1}}$ est défini comme le centralisateur}

La formule (96) implique que $T_G$ normalise $T_H$ (puisque $H$ est distingué dans $G$), donc comme $T_G/\mathfrak{z}$ est commutatif et $T_H\cap\mathfrak{z}\subset\unsure{\{H\cap\mathfrak{z}\}|H\cap\mathfrak{z}}=\{1\}$, on voit

\[T_G\text{ centralise }T_{\deleted{\unsure{G}}H}\tag*{(97)}\]

% archive sheet number: 249
\scanpage{251}

\mspage{128}

(alors qu'il n'est pas sûr que \deleted{\unsure{$B$}} $T_G$ soit comm., car il n'est pas sûr que $\unsure{\mathfrak{z}|\mathfrak{Z}}$ le soit \ldots). On voit aussi immédiatement

\[\left\{\begin{aligned}
&T_{\deleted{\unsure{H}}G}\text{ normalise }U_0\text{ et }U_1,\quad \inserted{\text{d'où}}\ T_G\cap H=T_H\\
&T_G\cap U_0=T_G\cap U_1=\{1\}\qquad(\text{car }T_H\cap U_0=T_H\cap U_1=\{1\})
\end{aligned}\right.\tag*{(98)}\]

\deleted{\unsure{on}} introduit donc les sous-groupes

\[\left.\begin{aligned}
B_{0G}^{\unsure{,|\prime}}&=T_G.U_0\\
B_{1G}^{\unsure{,|\prime}}&=T_G.U_1
\end{aligned}\right\}\text{ produits semi-directs}\tag*{(99)}\]

\unsure{donc|d'où}

\[\sigma(B_{0G})=B_{1G},\qquad \sigma(B_{1G})=B_{0G}.\tag*{(100)}\]

On trouve

\[\left\{\begin{aligned}
B_{0G}&=G_{D_0}=\{g\in G\mid gs_0\in Zs_0\}\\
B_{1G}&=G_{D_1}=\{g\in G\mid gt_0\in Zt_0\}
\end{aligned}\right.\tag*{(101)}\]

d'où

\[B_{0G}\cap B_{1G}=T_G.\tag*{(102)}\]

Pour prouver (101), il suffit de prouver que \inserted{si} $g\in G$, \deleted{$gs_0$} $\lambda\in Z$ sont tels que

\[gs_0=\lambda s_0\]

alors $g\in T_G U_0$. Or quitte à corriger \inserted{$g$} par un élément convenable de $T_G$, OPS $\delta(g)=1$ i.e. $g\in G_0$. \deleted{Quitte à \unsure{le}} \inserted{et quitte} à corriger par un élément de $\mathfrak{z}\subset T_G$, OPS $g\in H$ (cf plus bas). Mais on a alors par (\unsure{83|84}) $g\in B_{0H}\overset{\mathrm{déf}.}{=}T_H.U_0\subset T_G.U_0$ cqfd.

On a utilisé en passant

\[G_0=\mathfrak{z}.H\qquad (1/2\text{ direct}),\tag*{(103)}\]

% archive sheet number: 250
\scanpage{252}

qui résulte aussitôt des définitions et du fait que $H$ est simpl. transitif sur $\unsure{A_0|A_{\mathrm{o}}|\vec{A}_0}$. Notons aussi

\[\begin{aligned}
U_0\times T_G\times U_1&\unsure{-|\longrightarrow}G\qquad\text{est injectif}\\
(u_0,t,u_1)&\longmapsto u_0tu_1
\end{aligned}\tag*{(104)}\]

en effet, comme $T_G$ normalise $U_0,U_1$, cela se réduit aux formules

\[T_G\cap U_0=T_G\cap U_1=\{1\}\quad(98)\quad\text{et}\]

\[\underset{B_{0G}}{\underline{U_0.T_G}}\cap U_1=\{1\}\tag*{(*)}\]

Mais on a

\[\left\{\begin{aligned}
B_{0G}\cap H&=B_{0H}\quad(=U_0.T_H)\\
B_{1G}\cap H&=B_{1H}\quad(=U_1.T_H)
\end{aligned}\right.\tag*{(105)}\]

(résulte immédiatement de $T_G\cap H=T_H$), donc on se \deleted{réduit} ramène à la formule analogue dans $H$

\[U_0.T_H\cap U_1=\{1\},\]

qui est satisfaite par hypothèse.

On a un peu \unsure{lambiné|bambiné|bousculé}, sans aller droit à la construction de (64) en termes des données $(H,U_0,\sigma,T_H)$ -- je vais y \unsure{arriver|revenir}, en introduisant

\[\left\{\begin{aligned}
B'_{0\deleted{G}}&=G_{s_0}=\{g\in G\mid gs_0=s_0\}\subset B_{0G}=G_{D_0}.\\
B'_{1\deleted{G}}&=G_{t_0}=\{g\in G\mid gt_0=t\}\subset B_{1G}=G_{D_1}.
\end{aligned}\right.\tag*{(106)}\]

On trouve en fait, plus précisément

\[\left\{\begin{aligned}
B'_{0\deleted{G}}&=T_{0\deleted{G}}.U_0\\
B'_{1\deleted{G}}&=T_{1\deleted{G}}.U_1
\end{aligned}\right.\qquad\text{produits semi-directs}\tag*{(107)}\]

% archive sheet number: 251
\scanpage{253}

\mspage{129}

et de même

\[G=T_0.H=T'_0.H\qquad(\text{produits }1/2\text{ directs}),\tag*{(108)}\]

pour prouver (108), notons d'abord

\[T_0\cap H\subset T_G\cap H\deleted{\unsure{\subseteq}}\overset{(96)}{=}T_H\quad\text{donc}\]

\diagram{\begin{tikzcd}[ampersand replacement=\&,column sep=small]
{T_0\cap H=T_0\cap T_H=\operatorname{Ker}\bigl(} \&[-1.6em] T_H \arrow[r,hook,"\unsure{k|\kappa}"] \arrow[rr,bend right=22,swap,"\unsure{k|\kappa}"] \& T_G \arrow[r,"\unsure{\delta_0'|\delta_0}"] \& Z \&[-1.6em] {\bigr)=\{1\};}
\end{tikzcd}}

on voit de $=$ que $T'_0\cap H=\{1\}$. Pour voir que tout $g\in G$ s'écrit sous la forme $t_0.u$ ($t_0\in T_0,u\in H$), ou ce qui revient au $=$, sous la forme $u.t_0$ (puisque $T_0$ normalise $H$), OPS, comme $H$ est transitif sur $S$, quitte \deleted{\unsure{mod.}} \unsure{à} corriger \inserted{$g$} par un $u\in H$, que l'on a $gs_0=s_0$, donc on est ramené à la formule (107). Quitte à corriger par un élément de $T_0$ (puisque $T_0\xrightarrow{\delta|T_0} Z$ surjectif) OPS $\delta(g)=1$ i.e. $g\in G_0\cap B'_0$ et \deleted{comme on a vu que} quitte à corriger par un $z\in\unsure{\mathfrak{z}|\mathfrak{Z}}\subset T_0$, OPS $g\in H\cap B'_0\deleted{\unsure{=}}$, et il suffit de \deleted{voir} \inserted{noter} que

\[\begin{aligned}
H\cap B'_0&=U_0,\\
H\cap B'_1&=U_1,
\end{aligned}\tag*{(109)}\]

\deleted{\unsure{q}} ce qui achève de prouver (107) (108). On prouve de $=$

\[\left\{\begin{aligned}
B_{0G}&=T_0.B_{0H}\deleted{\unsure{=}}\\
B_{1G}&=T_{1}.B_{1H}
\end{aligned}\right.\ ,\tag*{(110)}\]

notons aussi

\[\left\{\begin{aligned}
\unsure{\hat B'_0|B'_0}\cap G_0&=\mathfrak{z}.U_0\\
B'_1\cap G_1&=\mathfrak{z}.U_1
\end{aligned}\right.\qquad\text{prod. }1/2\text{ directs}\tag*{(111)}\]

% archive sheet number: 252
\scanpage{254}

et

\[\left\{\begin{aligned}
B_{0G}\cap G_0&=\unsure{\mathfrak{z}|\mathfrak{Z}}.B_{0H}\\
B_{1G}\cap G_0&=\mathfrak{z}.B_{1H}
\end{aligned}\right.\qquad\text{produits semi-directs},\tag*{(\unsure{112|111})}\]

formules qui proviennent de la conjonction de (108), et du fait que par définition

\[G_0=\operatorname{Ker}(G\xrightarrow{\delta}Z)\]

et qu'enfin $\delta$ s'explicite comme le composé

\[G=T_i.H\xrightarrow{\mathrm{can}}T_i\xrightarrow{\delta|T_i}Z\qquad i\in\{0,1\}\]

et \unsure{qu'|qui|que} \deleted{\unsure{en}}enfin $\deleted{\unsure{T_i\cap T_H}}\mathfrak{z}$ est le noyau de $\delta|T_i$.

Une partie des relations explicitées dans les formules précédentes sont résumées dans le diagramme d'inclusions

\diagram{\resizebox{\linewidth}{!}{\begin{tikzcd}[ampersand replacement=\&,column sep=small,row sep=large]
\& U_0 \arrow[r,hook,no head,"Z"] \arrow[d,hook,swap,"\mathfrak{z}"] \& B_{0H}=T_H.U_0 \arrow[r,hook] \arrow[d,hook,swap,"\mathfrak{z}"] \& H \arrow[d,hook]\\
(113) \& \widetilde U_0=\mathfrak{z}.U_0 \arrow[r,hook,no head,"Z"] \arrow[d,hook,swap,"\widetilde Z"] \& \widetilde B_0=\mathfrak{z}.B_{0H}=(\mathfrak{z}\times T_H).U_0 \arrow[r,hook] \arrow[d,hook,"\widetilde Z"] \& G_0=\mathfrak{z}.H \arrow[d,hook,"\widetilde Z"]\\
\& B'_0=T_0.U_0 \arrow[r,hook,no head,"Z"] \& B_{0G}=T_0.B_{0H}=(T_0\times T_H).U_0 \arrow[r,hook] \& G=T_0.H
\end{tikzcd}}}

dont \deleted{\unsure{les}} quatre carrés composants sont cartésiens et cocartésiens, et où on a écrit, au dessus des inclusions distinguées, la valeur à isom près du quotient correspondant. J'ai posé

\[\widetilde Z\overset{\text{déf}}{=}G/H\underset{\mathrm{surj}.}{\overset{\unsure{c|e}}{\longrightarrow}}G/G_0\deleted{\unsure{\simeq}}\lhook\joinrel\xrightarrow{\delta} Z\tag*{(114)}\]

% margin note written sideways
\margin{on désigne par $\unsure{\widetilde\delta|\delta}:G\longrightarrow G/H\simeq\widetilde Z$ l'appl. can. de passage au quotient.}

d'où un hom can. composé \deleted{\unsure{de}} de noyau $\mathfrak{z}$

\[0\longrightarrow\mathfrak{z}\longrightarrow\widetilde Z\xrightarrow{c}Z\tag*{(115)}\]

% archive sheet number: 253
\scanpage{255}

\mspage{130}

qui \inserted{$(\unsure{c})$} est surjectif ssi on est dans la ``situation d'isomorphie'' i.e. $\delta:G\unsure{\longrightarrow|\longleftrightarrow} Z$ surjectif.

% margin note written sideways
\margin{\uline{NB} On a des iso canoniques $(115')\left\{\begin{aligned}T_0&\xrightarrow[\sim]{\unsure{\kappa_0|\pi_0}}\widetilde Z\\T_1&\xrightarrow[\sim]{\unsure{\kappa_1|\pi_1}}\widetilde Z\end{aligned}\right.$ mais on se gardera d'identifier $T_0,T_1$ (sous-groupes distincts de $G$) à $\widetilde Z$ (groupe quotient de $G$)}

On a un diagramme symétrique à (113) en y remplaçant l'indice $0$ par $1$, il est déduit de (113) en \deleted{y \unsure{remplaçant|transformé}} \inserted{le transformant} par l'autom $\operatorname{int}(\sigma)$, où $\sigma$ est considéré comme élément de $G$.

\deleted{Si on veut une description ``symétrique'' de $G$, \unsure{on} on peut écrire}

\[G=\deleted{\unsure{B}}T_G.H,\quad\text{avec }T_G\cap H=T_H\quad(98)\tag*{(116)}\]

donc dans cette description \inserted{de $G$}, il s'agit de bien expliciter

\[\left\{\begin{aligned}
&\text{a) le groupe }T_G,\quad\text{b) le plongement de }T_H\text{ dans }T_G,\text{ et}\\
&\text{\deleted{enfin} c) l'action de }T_G\text{ sur }H,\quad\text{\deleted{enfin} ce qui}\\
&\text{permet de reconstituer l'inclusion}\\
&\hspace{8em}H\subset G\\
&\text{de (64), il faut de plus}\\
&\text{d) expliciter l'hom d'inclusion } Z\subset T_G,\text{ et enfin}\\
&\text{e) expliciter l'hom }\delta|T_G:T_G\longrightarrow Z,\text{ pour}\\
&\text{réconstituer le diagramme (64).}
\end{aligned}\right.\tag*{(116')}\]

Nous allons expliciter $T_G$ \deleted{\unsure{comme} à l'aide de la suite exacte} comme coextension centrale

\[\deleted{1\longrightarrow\unsure{\mathfrak{z}|\mathfrak{Z}}\longrightarrow
\underbrace{T_{0\deleted{\unsure{G}}}\times T_1}_{\simeq\,\widetilde Z\times\widetilde Z}
\longrightarrow T_G\longrightarrow 1,}\tag*{\deleted{(117)}}\]

\deleted{en supposant pour simplifier que \unsure{$\widetilde Z_0$|$\widetilde T_0$} est contenu \inserted{(donc $T_0,T_1$)} dans $T_1$ aussi commutatifs, donc aussi \unsure{$Z$}} \inserted{\deleted{$T_0$ et $T_1$ isom à $\widetilde Z$ sont}}

\[\deleted{T_G\simeq T_H\times T_0\simeq Z\times\widetilde Z}\]

% archive sheet number: 254
\scanpage{256}

\[1\longrightarrow\underset{\substack{\shortparallel\\Z}}{T_H}\longrightarrow T_G\xrightarrow{\widetilde\delta|T_G}\widetilde Z\longrightarrow1\tag*{(117)}\]

qui est scindée par $T_{0G}$ et $T_{1G}$, qui donnent deux \deleted{\unsure{iso}} hom

\diagram{\begin{tikzcd}[ampersand replacement=\&]
(118)\quad\widetilde Z \arrow[r,shift left=2.6pt,"{\deleted{\unsure{\delta|T_0}}u_0}"] \arrow[r,shift right=2.6pt,swap,"u_1"] \& T_G
\end{tikzcd}}

satisfaisant la relation

\[\begin{aligned}
\underset{\substack{s_0\mapsto\deleted{\unsure{\widetilde{\phantom{\lambda}}}}\lambda s_0\\t_0\mapsto t_0}}{u_1(\widetilde\lambda)}\deleted{\unsure{u_0}}
={}&\underbrace{u_0(\widetilde\lambda)}_{\substack{(s_0\mapsto s_0)\\(t_0\unsure{-|\mapsto}\lambda t_0)}}
\deleted{\unsure{u_{c(\widetilde\lambda)}}}
\underbrace{\overbrace{u_{c(\widetilde\lambda)}}^{\text{él. central de }T_G}}_{\substack{s_0\mapsto\lambda s_0\\t_0\mapsto\lambda^{-1}t_0}}
\end{aligned}\tag*{(119)}\]

qui \deleted{\unsure{se}} s'écrit aussi

\[\begin{aligned}
u_1(\widetilde\lambda)u_0(\widetilde\lambda)^{-1}&=u_{c(\widetilde\lambda)}\\
\text{ou }\ u_0(\widetilde\lambda)^{-1}u_1(\widetilde\lambda)&
\end{aligned}\]

et comme les \deleted{\unsure{premiers}} deux mbes sont dans $H$, il suffit de vérifier que les valeurs dans $s_0,t_0$ sont les $=$s, ce qui est immédiat.

L'un et l'autre scindage de (117) donnent deux isomorphismes

\diagram{\begin{tikzcd}[ampersand replacement=\&]
(120)\quad {\deleted{\unsure{Z}}T_G} \arrow[r,shift left=2.5pt,"\sim"] \arrow[r,shift right=2.5pt,"\sim" description] \& Z\times\widetilde Z,
\end{tikzcd}}

le passage de l'un à l'autre se faisant via l'hom $\widetilde Z\xrightarrow{k}Z$.

Quand $\widetilde Z$ est commutatif, \inserted{donc $T_G$ commutatif,} cette description se simplifie, en écrivant la suite exacte

\[\begin{aligned}
0\longrightarrow\unsure{\mathfrak{z}|\mathfrak{Z}}\longrightarrow
&\overbrace{T_0\unsure{\mathbin{.}\times|\times}T_1}^{\simeq\widetilde Z\times\widetilde Z}\longrightarrow T_G\longrightarrow1\\
&(t,t')\longmapsto t.t'\\[-2pt]
z\longmapsto (z,z^{-1})&
\end{aligned}\tag*{(121)}\]

qui explicite $T_G$ comme un quotient de $T_0\times T_1\simeq\widetilde Z\times\widetilde Z$

% archive sheet number: 255
\scanpage{257}

\mspage{131}

\deleted{\unsure{dans} Dans $\widetilde Z\times\widetilde Z$ on a le sous-groupe diagonal, contenant $\unsure{\mathfrak{z}|\mathfrak{Z}}$, et le quotient de celui-ci par $\mathfrak{z}$ s'identifie alors au sous-groupe $Z\hookrightarrow T_G$, en d'autres termes, on a On retrouve le sous}

La suite exacte (121) se réécrit en termes de $\widetilde T$, et \unsure{reçoit|aussi|inscrit} la suite (115) comme indiqué dans le diagramme

\diagram{\resizebox{\linewidth}{!}{\begin{tikzcd}[ampersand replacement=\&,column sep=small,row sep=large]
\& 0 \arrow[r] \& \mathfrak{z} \arrow[r] \arrow[d,equal] \& \widetilde Z \arrow[r,"c",no head] \arrow[d,hook,"{\widetilde\lambda\mapsto(\widetilde\lambda,\widetilde\lambda^{-1})}"] \& Z \arrow[d,hook,"{\substack{\mathrm{incl}.\\\unsure{ca|\alpha}}}"] \& \\
(122) \& 0 \arrow[r] \& \mathfrak{z}\arrow[r,swap,"{z\mapsto(z,z^{-1})}"] \& \widetilde Z\times\widetilde Z \arrow[r,swap,"{(\widetilde\lambda,\widetilde\mu)\mapsto u_0(\widetilde\lambda)u_1(\widetilde\mu)}"] \& T_G\arrow[r] \&1
\end{tikzcd}}}

NB si $z\in\mathfrak{z}$, on a $u_0(z)=u_1(z)$

qui permet de retrouver l'inclusion canonique $Z\subset T_G$, quand $c:\widetilde Z\longrightarrow Z$ est surjective.

Ainsi, le groupe $T_G$ et l'inclusion $T_H\subset T_G$ \unsure{c.à.d} les données a)b) de (116'), s'explicitent en termes du groupe $\widetilde Z$ et \unsure{de|du} \inserted{l'hom $c:\widetilde Z\unsure{-|\longrightarrow}Z$ du} diagramme (115), qui fournit aussi l'hom. $\delta:T_G\unsure{-|\longrightarrow}Z$ \inserted{(donnée e) de (116'))} par commutativité, soit comme le composé de $\widetilde\delta|T_G:T_G\longrightarrow\widetilde Z$ et de $c:\widetilde Z\unsure{-|\longrightarrow} Z$, dans la \deleted{\unsure{suite}} description (117), soit (dans le cas de la description (122)), par commutativité du diagramme

% archive sheet number: 256
\scanpage{258}

\diagram{\begin{tikzcd}[ampersand replacement=\&,column sep=large,row sep=large]
\& \widetilde Z\times\widetilde Z \arrow[r] \arrow[d,"{(\widetilde\lambda,\widetilde\mu)\mapsto\widetilde\lambda+\widetilde\mu}"] \& T_G\arrow[r]\arrow[d,"\delta"] \& 1\\
(123) \& \widetilde Z\arrow[r,"c"] \& Z. \&
\end{tikzcd}}

Reste à décrire, \unsure{à} part \deleted{\unsure{l'incl|l'action}} (115) lui-$=$, et l'action de $T_G$ sur $H$ \inserted{(donnée c))}, i.e. de $T_0\xrightarrow{\sim}\widetilde Z$ et de $T_1\simeq\widetilde Z$ sur $H$, l'inclusion de $Z$ dans $T_G$ (donnée d) de (116')). \deleted{Je dis que $\unsure{c}$} \inserted{Cette inclusion} peut se décrire à l'aide de l'hom

\diagram{\begin{tikzcd}[ampersand replacement=\&]
(12{\deleted{\unsure{6}}4})\qquad c':Z \arrow[rr] \arrow[dr,swap,"{\deleted{\unsure{\mathrm{incl}.}}}","{\diagup}" description] \& \& \widetilde Z,\\
\& {\deleted{\unsure{T_G}}}\arrow[ur,"{\deleted{\unsure{\widetilde\delta}}}","{\diagup}" description] \&
\end{tikzcd}}

défini comme composé

\[Z\overset{\mathrm{incl}}{\unsure{-|\longrightarrow}}\deleted{\unsure{T_G}}G\xrightarrow{\widetilde\delta}\widetilde Z\quad\text{ou}\quad Z\overset{\mathrm{incl}}{\unsure{-|\longrightarrow}}T_G\xrightarrow{\widetilde\delta|T_G}\widetilde Z,\tag*{(12\unsure{5|4})}\]

qui satisfait la relation

\diagram{\begin{tikzcd}[ampersand replacement=\&,column sep=large]
(12\unsure{6|5})\quad Z \arrow[r,"c'"] \arrow[rr,bend right=24,"{\lambda\mapsto\lambda^2}" description] \& \widetilde Z\arrow[r,"c"] \& Z\qquad \begin{aligned}&\text{i.e.}\\&c(c'(\lambda))=\lambda^2\\&\forall\lambda\in Z.\end{aligned}
\end{tikzcd}}

Ainsi, quand on écrit

\[T_G=\deleted{Z}T_H\times\deleted{\widetilde Z}T_0\tag*{(12\unsure{6|5|7})}\]

la donnée de l'inclusion $Z\subset T_G$ revient à celle d'un hom

\[Z\hookrightarrow\underset{Z}{\underline{T_H}}\times\underset{\widetilde Z}{\underline{T_0}},\tag*{(12\unsure{7|6})}\]

donné par la compatibilité

% archive sheet number: 257
\scanpage{259}

\mspage{132}

\diagram{\resizebox{\linewidth}{!}{\begin{tikzcd}[ampersand replacement=\&,column sep=large,row sep=large]
(128)\quad Z \arrow[r,hook] \arrow[dr,swap,"{\lambda\mapsto(\lambda,c'(\lambda)\deleted{\unsure{^{-1}}})}"] \& T_H\unsure{\times|*}T_0\simeq T_H\times T_0 \arrow[d,leftrightarrow,"{\unsure{\sim|2|\wr}}"] \& \\
\& Z\times\widetilde Z \& \begin{aligned}(\lambda,\widetilde\mu)&\mapsto\deleted{\unsure{(}}u(\lambda).u_0(\widetilde\mu)\\ \text{i.e. }&\boxed{\lambda=u(\lambda)u_0(c'(\lambda))}\end{aligned}
\end{tikzcd}}}

Pour vérifier la commutativité de ce diagramme il suffit de \deleted{vérifier que} noter que les \deleted{\unsure{dif}} deux flèches \deleted{\unsure{envi}} $Z\longrightarrow T_H\times T_0$ de \unsure{(121)|(128)|(123)} ont $=$ composé avec $\underset{Z\times\widetilde Z}{\underline{T_H\times T_0}}\xrightarrow{\widetilde\delta}\widetilde Z$, donc les deux images d'un $\lambda$ sont congrues mod $T_H$, donc pour voir qu'elles sont égales, il suffit de voir que les \deleted{\unsure{composés|compositions}} actions sur $s_0,t_0$ sont égales i.e. égales à $\lambda s_0,\lambda t_0$, ce qui est immédiat. Si on avait utilisé l'isom $T_G\simeq T_H\times T_1$, et l'isom avec $Z\times\widetilde Z$ donné par

\[(\lambda,\widetilde\mu)\longmapsto\deleted{\unsure{\{} }u(\lambda).u_1(\widetilde\mu)\deleted{=u(\lambda).u_0(\widetilde\mu)u_{(c(\widetilde\lambda))}}\]

on trouve

\[\begin{aligned}
\lambda&\underset{(\unsure{128|123})}{=}u(\lambda).u_0(c'(\lambda))
\underset{(119)}{=}u(\lambda)u_1(c'(\lambda))u(\overbrace{c(c'(\lambda))}^{\lambda^2\text{ par }(126)})^{-1}\\
&=u(\lambda\deleted{\unsure{c(c'\lambda)}}\lambda^{-2})u_1(c'(\lambda))
\end{aligned}\]

d'où

\[\deleted{\unsure{\lambda=u(\lambda)^{-1}u_0(\lambda)^{-1}}}
\inserted{\lambda=u(\lambda)u_0(c'(\lambda))}
=u(\lambda^{-1})u_1(c'(\lambda))\qquad(\lambda\in Z).\tag*{(129)}\]

Dans le cas où $\widetilde Z$ est commutatif, et où $T_G$ est décrit par (121) i.e. la deuxième ligne de (122), on aimerait pouvoir interpréter

% archive sheet number: 258
\scanpage{260}

le sous-groupe $Z$ de $T_G$ comme une sorte de sous-groupe diagonal, avec un diagramme commutatif

\diagram{\begin{tikzcd}[ampersand replacement=\&,column sep=small]
\& 1\arrow[r] \& \unsure{\mathfrak{z}|\mathfrak{Z}}\arrow[r]\arrow[d,equal] \& \widetilde Z\arrow[r,"c"]\arrow[d,"\mathrm{diag}"] \& Z\arrow[r,""{name=nblink,coordinate}]\arrow[d,hook,"\mathrm{incl}."] \& 1 \& |[alias=nbtip]| {} \arrow[from=nblink,to=nbtip,bend left=43]\\
(130)\&1\arrow[r] \& \mathfrak{z}\arrow[r,swap,"{z\mapsto(z,z^{-1})}"] \& \widetilde Z\times\widetilde Z\arrow[r] \& T_G\arrow[r] \&1
\end{tikzcd}}

\margin{NB on a supposé $c$ surjectif (cas où la ``condition d'isomorphie'' est vérifiée)}

\deleted{déjà le} \deleted{mais ce diagramme n'est commutatif,}

l'inclusion $Z\subset T_G$ étant donc déduite (ou mieux sur le sous-groupe $c(\widetilde Z)$ de $Z$) par passage au quotient à partir de l'inclusion diagonale $\widetilde Z\longrightarrow\widetilde Z\times\widetilde Z$. Mais pour que cette inclusion \deleted{\unsure{transforme}} \inserted{applique bien} $\mathfrak{z}=\operatorname{Ker} c$ dans $\mathfrak{z}$ \deleted{$\mathfrak{z}=$} $\operatorname{Ker}(\widetilde Z\times\widetilde Z\longrightarrow T_G)$, il faut (et suffit) qu'on ait

\[z^2=1\qquad\forall z\in\mathfrak{z},\tag*{(131)}\]

et alors le premier carré de (130) est bien commutatif (l'hom $\mathfrak{z}\longrightarrow\widetilde Z\times\widetilde Z$ de la deuxième ligne \deleted{\unsure{s'identifie à}} s'écrit aussi, alors $z\unsure{-|\mapsto}(z,z)$ i.e. est \deleted{\unsure{induit par}} \inserted{le composé}

\[\mathfrak{z}\overset{\mathrm{diag}}{\unsure{-|\longrightarrow}}\mathfrak{z}\times\mathfrak{z}
\underset{\substack{\mathrm{incl}.\\\mathrm{can}}}{\lhook\joinrel\longrightarrow}
\widetilde Z\times\widetilde Z\quad).\]

La commutativité du deuxième carré de (130) s'écrit sous la forme

\[c(\widetilde\lambda)=u_0(\widetilde\lambda)u_1(\widetilde\lambda)\qquad\forall\widetilde\lambda\in\widetilde Z\tag*{(132)}\]

(qui \deleted{\unsure{imp}} pour $\deleted{\unsure{z}}\widetilde\lambda=z\in\mathfrak{z}$ se réduit à (131)). On voit immédiatement que les deux mbes de (132) \deleted{\unsure{appliq}} ont $=$ effet sur $(s_0,t_0)$

\[(s_0,t_0)\longmapsto(\lambda s_0,\lambda t_0),\quad\text{où }\lambda=c(\widetilde\lambda),\]

% archive sheet number: 259
\scanpage{261}

\mspage{133}

donc ils sont congrus dans $T_G$ modulo $\unsure{\mathfrak{z}|\mathfrak{Z}}$. Pour \deleted{voir} qu'ils soient égaux, il \deleted{suffit de} \inserted{\unsure{f. et s.|suffit}} que leurs images dans $G/H\simeq\widetilde Z$ soient les $=$s, puisque l'on a $\mathfrak{z}\cap H=\{1\}$, i.e. que leurs images par $\widetilde\delta$ soient les ms, ce qui s'écrit sous la forme

\[c'(c(\widetilde\lambda))=\widetilde\lambda^2\qquad\forall\widetilde\lambda\in\widetilde Z\tag*{(133)}\]

(à rapprocher de (127), qui implique que les deux mbes de (\unsure{133|138}) sont égaux mod $\mathfrak{z}$).

En résumé

\[\begin{aligned}
&\forall\widetilde\lambda\in\widetilde Z,\ \text{on a}\\
&\hspace{3em}\underset{\substack{\rotatebox{90}{$\in$}\\Z\subset T_G}}{c(\widetilde\lambda)}
\equiv u_0(\widetilde\lambda)u_1(\widetilde\lambda)\quad\bmod\mathfrak{z}.
\end{aligned}\tag*{(134)}\]

Pour que cette congruence soit une égalité, il f. et s. que l'on ait (133), (qui implique (131)).

Quand de plus \deleted{\unsure{si}} $c:\widetilde Z\unsure{-|\longrightarrow}Z$ est surjectif, la formule (133) redonne une façon d'expliciter le plongement de $Z$ dans $T_G$, en termes de la seule donnée du diagramme (126).

En résumé, pour reconstituer tout le diagramme (64) à isom près -- en même temps que (113), -- à partir des données (\unsure{89|88})

% archive sheet number: 260
\scanpage{262}

i.e. $(H,U_0\unsure{\in|\subset}H,\deleted{\unsure{\sigma\in H}}, T_H\subset H,\sigma\in H)$ satisfaisant les conditions

% slanted
\margin{$U_0,T_H$ ss-groupes de $H$}

\[\left\{\begin{aligned}
&\sigma^2=1\\
&U_0\cap U_1=\{1\}\qquad\text{où }U_1=\sigma(U_0)\\
&\sigma(t)=t^{-1}\qquad\forall t\in T_H\\
&T_H\text{ normalise }U_0\\
&U_0\times T_H\times U_1\unsure{-|\longrightarrow} H\qquad\text{injectif}\qquad(\text{où }U_1=\sigma(U_0)\\
&(u_0,t,u_1)\longmapsto u_0tu_1\hspace{12em}-
\end{aligned}\right.\tag*{(135)}\]

il reste à expliciter, en termes de ces données, le groupe $\widetilde Z$, les deux homomorphismes

\diagram{\begin{tikzcd}[ampersand replacement=\&]
{\deleted{\unsure{c,c'}}}\quad\widetilde Z\arrow[r,shift left,"c"] \& Z\arrow[l,shift left,"c'"]{\deleted{\unsure{=T_H}}}
\end{tikzcd}}

(où $Z$ est une copie conforme de $T_H$ via l'isom $u:Z\xrightarrow{\sim}T_H$)

et enfin l'opération de $\widetilde Z$ sur $H$, via

% slanted
\margin{donc c'est en fait l'opération de $T_0$ sur $H$}

\[\widetilde\lambda\longmapsto\operatorname{int}(u_0(\widetilde\lambda))|H\quad\deleted{\text{\unsure{écrire}}}\overset{\mathrm{déf}}{=}\Theta_0(\widetilde\lambda)\]

ou, au choix, via

% slanted
\margin{i.e. l'opération de $T_1$ sur $H$.}

\[\widetilde\lambda\longmapsto\operatorname{int}(u_1(\widetilde\lambda))|H\overset{\mathrm{déf}}{=}\Theta_1(\widetilde\lambda)\]

ces deux opérations étant \inserted{en effet} reliées par la formule -- déduite de (119)

\[u_1(\widetilde\lambda)=u_0(\widetilde\lambda)u(c(\widetilde\lambda))\]

qui donne

\[\Theta_1(\widetilde\lambda)=\Theta_0(\widetilde\lambda)\operatorname{int}(u(c(\widetilde\lambda))),-\tag*{(136)}\]

\deleted{Cela} de sorte que $\Theta_0(\widetilde\lambda),\Theta_1(\widetilde\lambda)$ se déterminent mutuellement. Il revient aussi au $=$ de \deleted{\unsure{de de|de}} demander, \inserted{plus symétriquement,} la construction du groupe $T_G$, comme groupe opérant sur $H$, avec dans $T_G$ les \deleted{deux \unsure{sous}-}\inserted{trois sous-}groupes $T_H\hookrightarrow T_G$ (plongement à déterminer), $T_0$ et $T_1$, \deleted{\unsure{ce qui est|que}} \inserted{$\widetilde Z$ étant} alors défini \deleted{\unsure{comme}} \inserted{de façon} \unsure{symétrique|symétriques} comme étant $T_G/T_H$, et l'isomorphisme $Z\xrightarrow{c'}\widetilde Z$ étant décrit comme le composé de

% archive sheet number: 261
\scanpage{263}

\mspage{134}

$Z\xrightarrow[\sim]{u}T_H\hookrightarrow T_G\xrightarrow{\mathrm{can}}\widetilde Z$. Il faudra de plus, \unsure{enfin|encore}, décrire l'hom. $c:\widetilde Z\longrightarrow Z$.

Rappelons-nous que l'on a

\[\left\{\begin{aligned}
T_0&=\left\{g\in\unsure{\mathfrak{S}_S|\mathscr{G}_S}\ \middle|\ 
\begin{aligned}&g\text{ normalise }H\subset\mathfrak S_{\unsure{S|z}}\\&g\text{ centralise } Z\subset\mathfrak S_S\\&g\deleted{\unsure{\in}}s_0=s_0\\&gt_0\in Z.t_0\end{aligned}\right\}\deleted{\Longleftrightarrow}\\
T_1&=\left\{g\in\mathfrak S_S\ \middle|\ 
\begin{aligned}&g\text{ normalise }H\\&g\text{ centralise }Z\\&gt_0=t_0\\&gs_0\in Z.s_0\end{aligned}\right\}
\end{aligned}\right.\ .\tag*{(137)}\]

\margin{$T_G=\left\{g\in\mathfrak S_S\ \middle|\ \begin{smallmatrix}g\text{ norm. }H\\g\text{ centr. }Z\\gs_0\in Zs_0\\gt_0\in Zt_0\end{smallmatrix}\right.$}

Or pour un espace homogène $S$ sous un groupe $H\inserted{\subset\mathfrak S_S}$, $S$ \unsure{ponctué|pointé} par une ``origine $s_0$'', les groupe des $g\in\mathfrak S_S$ qui normalisent $H$ et fixent $s_0$, i.e. des autom. de la situation

\[(S,H\subset\mathfrak S_S,s_0\in S)\longleftrightarrow(H,U_0=H_{s_0}),\quad\text{est can.}\]

isomorphe au groupe des automorphismes \inserted{$\bar g\deleted{\unsure{g}}$} de $H$ qui \deleted{\unsure{stabilisent}} normalisent $U_0$. \deleted{\unsure{Le ca}} Pour un \deleted{tel} $g$, dire qu'il \deleted{centralise $Z\subset\mathfrak S_S$, défini} \inserted{normalise un sous-groupe} $Z\subset\mathfrak S_S$ commutant à \deleted{\unsure{$G$}} $H$, défini par un sous-groupe $\underset{\deleted{\unsure{T_H}}}{\underline{B_{0H}}}\supset U_0$ du normalisateur de $U_0$, \unsure{se} signifie que l'autom. correspondant $\bar g$ de $H$ normalise $B_{0H}$. -- et l'autom \deleted{\unsure{induit sur}} $\inserted{\operatorname{int}(g)|Z}$ s'identifie alors à $\operatorname{int}(\bar g)|(B_{0H}/U_0)$.

Donc dire que $g$ \uline{centralise} $Z$, signifie que $\bar g$ normalise $B_{0H}$ et induit l'identité sur $B_{0H}/U_0$.

Il reste à exprimer la condition $gt_0\in Z.t_0$ en termes de $\bar g$, or \deleted{$gt_0$} $t_0$ correspond à $\sigma.U_0$, $gt_0$ à $\bar g(\sigma).U_0$, et \deleted{on voit qu'il existe} \inserted{$\lambda t_0$ ($\lambda\in Z$) \unsure{corres}}

% archive sheet number: 262
\scanpage{264}

\deleted{$\lambda\in T_H$ tel que $\unsure{\lambda u(\sigma)|\widetilde\lambda u(\sigma)}U_0=$}

pond à $\sigma.u(\lambda)U_0$, donc la condition

\[gt_0=\lambda t_0\]

s'écrit aussi

\[\bar g(\sigma)U_0=\sigma.u(\lambda)U_0\qquad\begin{aligned}&\text{pour }\lambda\in Z\text{ convenable}\\&\text{(d'ailleurs unique)}\end{aligned}\tag*{(13\unsure{8|6})}\]

En fait, ce\unsure{tte|s} conditions impliquent déjà la condition, en apparence plus forte, que $\bar g$ \inserted{(normalise $U_0$ et)} \uline{centralise} $T_H$ lui-$=$, et que de plus on a

\[\bar g(\sigma)=\sigma u(\lambda)\qquad\text{pour }\lambda\in Z\text{ convenable},\tag*{(13\unsure{9|8})}\]

soit encore

\[\bar g(\sigma)\in\sigma.T_H,\tag*{(1\unsure{40|39})}\]

\deleted{puisque} \inserted{En effet,} en termes de $g\in\inserted{T_0\subset}G\subset\unsure{\mathfrak S_S|\mathscr G_S}$, \deleted{cette relation} \inserted{on sait} bien que $g$ centralise $T_H$, et (139) s'écrit \deleted{s'écrit}

\[g\sigma g^{-1}=\sigma u(\lambda)\qquad\text{i.e.}\qquad\sigma(g)g^{-1}=u(\lambda)\]

ou encore, \deleted{en écrivant $t$ au lieu de $g$ (puisque $g\in T_0$)}

\[\overbrace{\sigma(\deleted{t}g)}^{\rotatebox{270}{$\in$}\, T_1}=u(\lambda)\overset{\substack{T_0\\\rotatebox{270}{$\in$}}}{\deleted{\unsure{t}}g}\qquad \deleted{\text{avec }\lambda=\delta(g)},\ \substack{\text{pour }g\in T_0\\\text{et }\lambda\in Z\text{ conv.}}\tag*{(141)}\]

\deleted{soit encore, posant $\widetilde\delta(t)=\widetilde\lambda$}

\[\deleted{u_1(\widetilde\lambda)=u(\lambda)u_0(\widetilde\lambda)}\]

\deleted{On a vu} \inserted{En effet,} les deux mbes sont dans $T_G$ (qui est normalisé par $\sigma$), et on a vu qu'ils sont congrus mod $U_0$, donc congrus mod $T_G\cap U_0=T_H\cap U_0=\{1\}$, donc ils sont égaux. De plus, \deleted{en prenant appliquant} \inserted{cette formule montre} que \inserted{$g$ et} $\sigma(g)$ \deleted{et $g$} ont $=$ image \inserted{$\widetilde\lambda$} dans $\widetilde Z=T_G/T_H$, donc, \deleted{comme $g$ \unsure{prouve}} \inserted{comme ils sont} dans $T_0$ et $T_1$ respectivement, ils peuvent s'écrire $u_0(\widetilde\lambda),u_1(\widetilde\lambda)$, et (141) s'écrit

\[u_1(\widetilde\lambda)=u(\lambda)u_0(\widetilde\lambda)\tag*{(142)}\]

qui montre qu'on a (cf (119))

\[\lambda=c(\widetilde\lambda).\tag*{(143)}\]

% archive sheet number: 263
